\begin{table*}[ht]
\centering
\caption{Positioning relative to representative misinformation simulation models. ``Empirical'' identifies inputs or validation data reported by each study; it does not imply predictive validity.}
\label{tab:related-models}
\small
\begin{tabular}{p{2.5cm}p{2.8cm}p{2.6cm}p{2.5cm}p{3.4cm}}
\toprule
Study & Main mechanism & Network/input & Intervention & Distinctive analytical focus \\
\midrule
\citet{sulis2020} & Competing misinformation and debunking states & Configurable synthetic networks & Fact-checking probability & Diffusion and forgetting \\
\citet{hills2018} & Evolutionary deception and cooperation & Random interactions; empirical pattern comparison & Behavioural costs & Deception types and cooperation \\
\citet{gausen2022} & Beliefs and newsfeed curation & Parameterisation and case comparisons using Twitter data & Curation objectives & Algorithmic misinformation--polarisation trade-off \\
\citet{keijzer2021} & Cultural influence by bots and humans & Synthetic network structures & Bot activity and connectivity & Direct versus indirect bot influence \\
\citet{puri2024} & Language-sensitive digital clones & Observed network and message histories & Counterfactual content & Reproduction of a particular online system \\
This study & Repeated source choice from accumulated evidence & Survey-informed source perceptions; synthetic directed networks & Source imitation, reach and recommendation policy & Adaptive source presentation with objective false-publication probability held separate \\
\bottomrule
\end{tabular}
\end{table*}
